%!TEX root = ../../../note

\subsection{Lecture 1}

\subsubsection{What Is Calculus?}
Differential and integral calculus are two branches of calculus.

\begin{description}
    \item[Differential calculus:] The study of slopes and tangent lines.
    \item[Integral calculus:] The study of areas and accumulation.
\end{description}

\subsubsection{Numbers}
The number systems discussed in this lecture include
\begin{align*}
    \N &= \set{1,2,3,\ldots}, \\
    \Z &= \set{\ldots,-1,0,1,\ldots}, \\
    \Q &= \set{\frac{m}{n} : m,n\in\Z,\ n\ne 0}.
\end{align*}

The natural numbers are closed under addition. The integers are closed under addition, subtraction, and multiplication, but they are not closed under division; for example, $1/2\notin\Z$. Therefore, we introduce the rational numbers to include quotients of integers with nonzero denominators.

\begin{theorem}[Irrationality of $\sqrt{2}$]
    The number $\sqrt{2}$ is irrational; equivalently, $\sqrt{2}\notin\Q$.
    \begin{proof}
        Suppose, for contradiction, that $\sqrt{2}\in\Q$. Then we can write
        \begin{equation*}
            \sqrt{2}=\frac{m}{n},
            \qquad m,n\in\Z,
            \qquad n>0,
            \qquad \gcd(m,n)=1.
        \end{equation*}
        Squaring both sides gives
        \begin{equation*}
            2=\frac{m^2}{n^2}
            \qquad\Longrightarrow\qquad
            m^2=2n^2.
        \end{equation*}

        Thus, $m^2$ is even, so $m$ is even. Write $m=2\ell$ for some $\ell\in\Z$. Substituting this into the equation above gives
        \begin{align*}
            2n^2 &= 4\ell^2, \\
            n^2 &= 2\ell^2.
        \end{align*}
        Therefore, $n^2$ is even, and hence $n$ is even as well.

        This means that both $m$ and $n$ are even, contradicting $\gcd(m,n)=1$. Therefore, $\sqrt{2}\notin\Q$.
    \end{proof}
\end{theorem}

The irrational numbers are the real numbers that are not rational:
\begin{equation*}
    \R\setminus\Q
    =\set{x\in\R : x\notin\Q}.
\end{equation*}

\begin{example}[Decimal representation]
Every real number has a decimal representation. In general, a decimal expansion can be written as
\begin{equation*}
    x=a+\sum_{i=1}^{\infty}\frac{x_i}{10^i},
    \qquad a\in\Z,
    \qquad x_i\in\set{0,1,2,\ldots,9}.
\end{equation*}
\end{example}
